(1) 中心力为 \(-GMm/r^2\)，角动量初始为零且守恒，因此轨道完全径向。总能量为释放时的势能 \(E=-GMm/R\)，有
\[
\frac m2\dot r^2-\frac{GMm}{r}=-\frac{GMm}{R},\quad
\dot r=-\sqrt{2GM\left(\frac1r-\frac1R\right)}.
\]
坠落取负根，因为半径随时间减小。

(2) 分离变量
\[
dt=-\frac{dr}{\sqrt{2GM(1/r-1/R)}}.
\]
令 \(r=R\cos^2\eta\)，其中从释放点起 \(\eta\) 从零增加。则 \(dr=-2R\sin\eta\cos\eta\,d\eta\)，根号中的因子为 \(\tan\eta/\sqrt R\)。于是
\[
t(r)=\sqrt{\frac{R^3}{2GM}}
\left[\eta+\sin\eta\cos\eta\right],
\qquad \eta=\arccos\sqrt{r/R}.
\]
恢复为 r 的显式形式：
\[
t(r)=\sqrt{\frac{R^3}{2GM}}\left[
\arccos\sqrt{\frac rR}+\sqrt{\frac rR(1-\frac rR)}\right].
\]
直接对时间表达式求 r 导数可恢复上述分离变量方程。

(3) 当 r→0⁺ 时 \(\eta\to\pi/2\)，所以理想点质量模型的到达中心时间
\[
t_{\rm center}=\frac\pi2\sqrt{\frac{R^3}{2GM}}.
\]
真实天体有限半径使粒子先碰撞表面，因此此式在有限半径天体中仅是数学外推而非真实撞击时间。

(4) 释放初期 \(r=R-\delta\)，\(\delta\ll R\)，能量关系给 \(\dot r^2\simeq 2GM\delta/R^2\)，积分得 \(\delta\simeq GMt^2/(2R^2)\)，正是初始加速度 \(GM/R^2\) 的常加速度位移。随着 r→0，\(|\dot r|\sim\sqrt{2GM/r}\) 发散，但剩余时间积分 \(\int_0^r\sqrt{s/(2GM)}\,ds\) 仍有限，且按 \(r^{3/2}\) 趋零。
